\documentclass[12pt,letterpaper,onecolumn]{exam}
\usepackage[utf8]{inputenc}
\usepackage[T1]{fontenc}
\usepackage{amsmath}
\usepackage{amssymb}
\usepackage[lmargin=71pt,tmargin=1.2in]{geometry}
\usepackage{xcolor}
\usepackage[brazilian]{babel}
\usepackage{natbib}
\usepackage{enumitem}
\usepackage{booktabs}
\usepackage[hidelinks]{hyperref}
\urlstyle{same}

\renewenvironment{solution}
{\par\noindent\textbf{R:}\enspace\color{blue}}
{\par\color{black}}

\thispagestyle{empty}

\begin{document}

\begingroup
    \centering
    \LARGE Lista 12 Gabarito - Econometria I - 2026\\[0.5em]
    \large Prof: André Luiz Pereira Mancha \par
    \large Monitor: Tuffy Licciardi Issa \par
\endgroup
\rule{\textwidth}{0.4pt}
\pointsdroppedatright
\printanswers
\renewcommand{\solutiontitle}{\noindent\textbf{Ans:}\enspace}
\newcommand{\qstar}{\hspace{-0.6em}\textsuperscript{*}\hspace{0.2em}}

\begin{questions}

\question\qstar (17.2, \citep{wooldridge2016}) Defina \textit{grad} como uma variável \textit{dummy} que indica se um estudante-atleta de uma grande universidade se formará dentro de cinco anos. Sejam \textit{hsGPA} a nota média no ensino médio e \textit{SAT} a pontuação total no exame SAT. Defina \textit{study} como o número médio de horas por semana passadas em uma sala de estudo organizada. Suponha que, usando os dados de 420 estudantes-atletas, obtenha-se o seguinte modelo logit:
\[
P(grad = 1 \mid hsGPA, SAT, study)
=
\Lambda(-1{,}17 + 0{,}24\,hsGPA + 0{,}00058\,SAT + 0{,}073\,study),
\]
em que \(\Lambda(z)=\exp(z)/[1+\exp(z)]\) é a função logit.
\begin{enumerate}[label=(\roman*)]
    \item Mantendo fixos \(hsGPA\) em \(3{,}0\) e \(SAT\) em \(1.200\), calcule a diferença estimada na probabilidade de formatura de alguém que tenha passado dez horas por semana em uma sala de estudo e de alguém que tenha passado cinco horas por semana.
    \item Explique.
\end{enumerate}

\begin{solution}
\begin{enumerate}[label=(\roman*)]
    \item Para \(study=10\),
    \[
    z_{10}=-1{,}17+0{,}24(3{,}0)+0{,}00058(1200)+0{,}073(10)=0{,}976.
    \]
    Logo,
    \[
    \Lambda(z_{10})=\frac{e^{0{,}976}}{1+e^{0{,}976}}\approx 0{,}726.
    \]
    Para \(study=5\),
    \[
    z_{5}=-1{,}17+0{,}24(3{,}0)+0{,}00058(1200)+0{,}073(5)=0{,}611,
    \]
    de modo que
    \[
    \Lambda(z_{5})=\frac{e^{0{,}611}}{1+e^{0{,}611}}\approx 0{,}648.
    \]
    Portanto, a diferença estimada é
    \[
    0{,}726-0{,}648=0{,}078.
    \]
    Assim, passar de cinco para dez horas semanais na sala de estudo eleva a probabilidade estimada de formatura em aproximadamente \(7{,}8\) pontos percentuais.

    \item O efeito não é exatamente igual a \(5\times 0{,}073\), porque \(0{,}073\) é o efeito sobre o índice latente do logit, e não diretamente sobre a probabilidade. A relação entre o índice e a probabilidade é não linear. Por isso, a variação na probabilidade deve ser calculada comparando \(\Lambda(z_{10})\) e \(\Lambda(z_{5})\).
\end{enumerate}
\end{solution}

\vspace{0.7cm}

\question (17.5, \citep{wooldridge2016}) Defina \textit{patents} como o número de patentes requeridas por uma firma durante determinado ano. Assuma que o valor esperado condicional de \textit{patents}, dados \textit{sales} e \(PD\), é
\[
E(patents \mid sales,PD)=\exp\{\beta_0+\beta_1\log(sales)+\beta_2 PD+\beta_3 PD^2\},
\]
em que \textit{sales} representa as vendas anuais da firma e \(PD\) é o total de gastos com pesquisa e desenvolvimento nos últimos três anos.
\begin{enumerate}[label=(\roman*)]
    \item Como você estimaria os \(\beta_j\)? Escreva sua resposta detalhando a natureza de \textit{patents}.
    \item Como você interpreta \(\beta_1\)?
    \item Encontre o efeito parcial de \(PD\) sobre \(E(patents \mid sales,PD)\).
\end{enumerate}

\begin{solution}
\begin{enumerate}[label=(\roman*)]
    \item Como \textit{patents} é uma variável de contagem não negativa, o procedimento natural é usar um modelo de regressão de Poisson por máxima verossimilhança, especificando
    \[
    E(patents \mid sales,PD)=\exp\{\beta_0+\beta_1\log(sales)+\beta_2 PD+\beta_3 PD^2\}.
    \]
    Mesmo que a distribuição condicional não seja exatamente Poisson, o estimador de quase-máxima verossimilhança de Poisson continua apropriado desde que a média condicional esteja corretamente especificada.

    \item Como \(\log(sales)\) entra linearmente no índice, \(\beta_1\) é uma elasticidade semielástica da média condicional: um aumento de \(1\%\) em \textit{sales} altera \(E(patents\mid sales,PD)\) em aproximadamente \(\beta_1/100\) em termos proporcionais. De forma equivalente, um aumento de \(1\%\) nas vendas altera o número esperado de patentes em aproximadamente \(100\beta_1/100=\beta_1\%\).

    \item Seja
    \[
    \mu(x)=E(patents \mid sales,PD)=\exp\{\beta_0+\beta_1\log(sales)+\beta_2 PD+\beta_3 PD^2\}.
    \]
    Então,
    \[
    \frac{\partial \mu(x)}{\partial PD}
    =
    \mu(x)\,(\beta_2+2\beta_3 PD).
    \]
    Portanto, o efeito parcial de \(PD\) é
    \[
    \frac{\partial E(patents \mid sales,PD)}{\partial PD}
    =
    E(patents \mid sales,PD)\,(\beta_2+2\beta_3 PD).
    \]
\end{enumerate}
\end{solution}

\vspace{0.7cm}

\question\qstar (17.6, \citep{wooldridge2016}) Considere uma função de poupança familiar para a população de todas as famílias dos Estados Unidos:
\[
sav = \beta_0 + \beta_1 inc + \beta_2 hsize + \beta_3 educ + \beta_4 age + u,
\]
em que \textit{hsize} é o tamanho da família, \textit{educ} são anos de escolaridade do chefe da família e \textit{age} é a idade do chefe da família. Assuma que
\[
E(u \mid inc,hsize,educ,age)=0.
\]
\begin{enumerate}[label=(\roman*)]
    \item Suponha que a amostra inclua apenas famílias cujo chefe tenha mais de 25 anos. Se usarmos o MQO nessa amostra, obteremos estimadores não viesados dos \(\beta_j\)? Explique.
\end{enumerate}

\begin{solution}
Sim, desde que a seleção da amostra dependa apenas das variáveis explicativas observadas, isto é, da restrição \(age>25\), e que a condição de média condicional nula continue válida na amostra selecionada:
\[
E(u \mid inc,hsize,educ,age,\ age>25)=0.
\]
Como a seleção é baseada em uma variável explicativa já presente no modelo, não há problema de viés de seleção por si só. O que muda é que a inferência passa a valer para a subpopulação com chefes de família acima de 25 anos. Portanto, o MQO continua produzindo estimadores não viesados e consistentes dos coeficientes, sob essa condição adicional.
\end{solution}

\vspace{0.7cm}

\question\qstar (17.8, \citep{wooldridge2016}) Considere o problema de avaliação de programas, em que \(y(0)\) e \(y(1)\) são os resultados potenciais sem e com tratamento, respectivamente. Seja \(w\) o indicador binário de tratamento. Sob as funções médias condicionais, defina
\[
m_0(x)=E[y(0)\mid x]
\qquad \text{e} \qquad
m_1(x)=E[y(1)\mid x].
\]
Lembre-se de que o efeito médio do tratamento é
\[
\tau_{ate}=\mu_1-\mu_0=E[y(1)]-E[y(0)]
\]
e que também pode ser escrito como
\[
\tau_{ate}=E[m_1(x)]-E[m_0(x)].
\]
\begin{enumerate}[label=(\roman*)]
    \item Explique por que
    \[
    \tau_{ate}=n^{-1}\sum_{i=1}^{n}\big[m_1(x_i)-m_0(x_i)\big]
    \]
    se \(x_1,\ldots,x_n\) forem tratados como uma amostra aleatória.
    \item Se você conhecesse as funções \(m_0(x)\) e \(m_1(x)\) e dispusesse de uma amostra aleatória \(x_1,\ldots,x_n\), explique por que
    \[
    \hat{\tau}_{ate}=n^{-1}\sum_{i=1}^{n}\big[m_1(x_i)-m_0(x_i)\big]
    \]
    seria um estimador não viesado de \(\tau_{ate}\).
    \item Agora considere os modelos \(m_0(x,\theta_0)\) e \(m_1(x,\theta_1)\), em que \(\theta_0\) e \(\theta_1\) são parâmetros e estimadores \(\hat{\theta}_0\) e \(\hat{\theta}_1\). Defina
    \[
    \hat{\tau}_{ate}=n^{-1}\sum_{i=1}^{n}\big[m_1(x_i,\hat{\theta}_1)-m_0(x_i,\hat{\theta}_0)\big].
    \]
    Como você estimaria \(\tau_{ate}\)?
    \item Defina a resposta observada como de costume:
    \[
    y=(1-w)y(0)+wy(1).
    \]
    Mostre que, se \(w\) é independente de \([y(0),y(1)]\) condicional em \(x\), então
    \[
    E(y\mid w,x)=(1-w)m_0(x)+wm_1(x).
    \]
    Portanto,
    \[
    E(y\mid w=0,x)=m_0(x)
    \qquad \text{e} \qquad
    E(y\mid w=1,x)=m_1(x).
    \]
    \item Se \(y(0)\) e \(y(1)\) seguirem modelos logit ou probit com probabilidades de resposta:
    \[
    G(\alpha_0+x'\beta_0)
    \qquad \text{e} \qquad
    G(\alpha_1+x'\beta_1),
    \]
    respectivamente, como você estimaria todos os parâmetros? Como você estimaria \(\tau_{ate}\)?
\end{enumerate}

\begin{solution}
\begin{enumerate}[label=(\roman*)]
    \item Como
    \[
    \tau_{ate}=E[m_1(x)-m_0(x)],
    \]
    e \(x_1,\ldots,x_n\) são uma amostra aleatória, a média amostral de \(m_1(x_i)-m_0(x_i)\) é a contraparte natural da esperança populacional. Assim,
    \[
    \tau_{ate}\approx n^{-1}\sum_{i=1}^{n}[m_1(x_i)-m_0(x_i)].
    \]
    Em amostras grandes, pela Lei dos Grandes Números, essa média converge para \(E[m_1(x)-m_0(x)]\).

    \item Como \(x_1,\ldots,x_n\) formam uma amostra aleatória,
    \[
    E\!\left[n^{-1}\sum_{i=1}^{n}\big(m_1(x_i)-m_0(x_i)\big)\right]
    =
    E[m_1(x)-m_0(x)]
    =
    \tau_{ate}.
    \]
    Portanto, o estimador proposto é não viesado para \(\tau_{ate}\).

    \item O procedimento é estimar \(\theta_0\) e \(\theta_1\) pelos métodos adequados ao modelo escolhido e, em seguida, substituir esses estimadores na expressão
    \[
    \hat{\tau}_{ate}=n^{-1}\sum_{i=1}^{n}[\hat{m}_1(x_i)-\hat{m}_0(x_i)].
    \]
    Ou seja,
    \[
    \hat{m}_0(x_i)=m_0(x_i,\hat{\theta}_0),
    \qquad
    \hat{m}_1(x_i)=m_1(x_i,\hat{\theta}_1),
    \]
    e calcula-se a média amostral das diferenças previstas.

    \item Como
    \[
    y=(1-w)y(0)+wy(1),
    \]
    temos
    \[
    E(y\mid w,x)=(1-w)E[y(0)\mid w,x]+wE[y(1)\mid w,x].
    \]
    Sob independência condicional de \(w\) em relação a \([y(0),y(1)]\), segue que
    \[
    E[y(0)\mid w,x]=E[y(0)\mid x]=m_0(x)
    \]
    e
    \[
    E[y(1)\mid w,x]=E[y(1)\mid x]=m_1(x).
    \]
    Logo,
    \[
    E(y\mid w,x)=(1-w)m_0(x)+wm_1(x).
    \]
    Portanto,
    \[
    E(y\mid w=0,x)=m_0(x)
    \qquad \text{e} \qquad
    E(y\mid w=1,x)=m_1(x).
    \]

    \item Esses resultados mostram que podemos estimar separadamente, nos grupos \(w=0\) e \(w=1\), dois modelos logit ou probit:
    \[
    \hat{m}_0(x)=G(\hat{\alpha}_0+x'\hat{\beta}_0),
    \qquad
    \hat{m}_1(x)=G(\hat{\alpha}_1+x'\hat{\beta}_1).
    \]
    Por fim, calcula-se
    \[
    \hat{\tau}_{ate}=n^{-1}\sum_{i=1}^{n}\big[\hat{m}_1(x_i)-\hat{m}_0(x_i)\big].
    \]
\end{enumerate}
\end{solution}

\vspace{0.7cm}

\question\qstar (17.C1, \citep{wooldridge2016}) Use os dados de \texttt{PNTSPRD} neste exercício. A variável \textit{favwin} é binária e vale um se o time favorito vence a partida. Um modelo de probabilidade linear para estimar a probabilidade de o time favorito vencer é:
\[
P(favwin=1\mid spread)=\beta_0+\beta_1 spread.
\]
\begin{enumerate}[label=(\roman*)]
    \item Explique por que, caso a aposta incorpore todas as informações relevantes, esperaríamos \(\beta_0=0{,}5\).
    \item Estime o modelo do item (i) por MQO. Teste \(H_0:\beta_0=0{,}5\) em relação a uma alternativa bilateral. Use erros padrão usuais e robustos em relação à heteroscedasticidade.
    \item Qual é a probabilidade estimada de o time favorito vencer quando \(spread=10\) e ela é estatisticamente significante?
    \item Agora, estime um modelo Probit para \(P(favwin=1\mid spread)\). Interprete e teste a hipótese nula de que o intercepto é zero. [Dica: lembre-se de que \(\Phi(0)=0{,}5\).]
    \item Use o modelo Probit para estimar a probabilidade de que o time favorito vença, caso \(spread=10\). Compare esse resultado com a estimativa MQO do item (iii).
    \item Adicione as variáveis \textit{favhome}, \textit{fav25} e \textit{und25} ao modelo Probit e teste a significância conjunta dessas variáveis usando o teste de razão de verossimilhança (quantos gl estão na distribuição qui-quadrado?). Interprete esse resultado, focando na possibilidade de que o \textit{spread} incorpore todas as informações observáveis antes de um jogo.
\end{enumerate}

\begin{solution}
\begin{enumerate}[label=(\roman*)]
    \item Se \(spread=0\), as casas de aposta consideram os times equilibrados. Nesse caso, a probabilidade prevista de vitória do favorito deve ser \(0{,}5\). Como, no modelo linear,
    \[
    P(favwin=1\mid spread=0)=\beta_0,
    \]
    segue que, sob incorporação completa das informações relevantes, esperaríamos \(\beta_0=0{,}5\).

    \item A parte empírica deve ser feita no software com MQO e, para o teste bilateral, usa-se
    \[
    t=\frac{\hat{\beta}_0-0{,}5}{se(\hat{\beta}_0)},
    \]
    tanto com erro padrão usual quanto com erro padrão robusto. A conclusão depende dos resultados numéricos da amostra.

    \item A probabilidade estimada pelo modelo linear é
    \[
    \widehat{P}(favwin=1\mid spread=10)=\hat{\beta}_0+10\hat{\beta}_1.
    \]
    Para verificar a significância, testa-se \(H_0:\beta_0+10\beta_1=0\) ou, mais relevantemente, examina-se se a probabilidade prevista é estimada com precisão, usando o erro padrão da combinação linear. Numericamente, isso deve ser calculado após a estimação.

    \item No probit,
    \[
    P(favwin=1\mid spread)=\Phi(\gamma_0+\gamma_1 spread).
    \]
    Se \(\gamma_0=0\), então
    \[
    P(favwin=1\mid spread=0)=\Phi(0)=0{,}5.
    \]
    Portanto, a hipótese \(H_0:\gamma_0=0\) verifica se, quando não há favoritismo pelo \textit{spread}, a probabilidade prevista de vitória é \(50\%\). O teste é um teste \(t\) ou \(z\) usual sobre o intercepto estimado.

    \item A probabilidade prevista no probit para \(spread=10\) é
    \[
    \widehat{P}(favwin=1\mid spread=10)=\Phi(\hat{\gamma}_0+10\hat{\gamma}_1).
    \]
    Essa previsão deve ser comparada com a do MQO,
    \[
    \hat{\beta}_0+10\hat{\beta}_1.
    \]
    Em geral, os valores são próximos no centro da distribuição, mas o probit garante previsões entre zero e um.

    \item O teste de razão de verossimilhança compara o modelo restrito, com apenas \textit{spread}, ao modelo irrestrito, com \textit{spread}, \textit{favhome}, \textit{fav25} e \textit{und25}. A estatística é
    \[
    LR=2(\ell_{irrestrito}-\ell_{restrito}),
    \]
    que, sob \(H_0\), tem distribuição qui-quadrado com \(3\) graus de liberdade. Se \(H_0\) não for rejeitada, isso sugere que, dado o \textit{spread}, essas variáveis adicionais não trazem informação relevante, reforçando a ideia de que o \textit{spread} já resume bem as informações observáveis antes do jogo. Se \(H_0\) for rejeitada, então o \textit{spread} não incorpora toda a informação disponível.
\end{enumerate}
\end{solution}

\end{questions}

\bibliographystyle{apalike}
\bibliography{refs}
\end{document}
